\begin{textsample}{POS Dim 3 – vem\_ed\_21 – Score 3.00 – vem\_ed\_21/t107.txt}  \label{ex:f3_pos_109}
Comunicações Do quebra-gelo à \textbf{institucionalização} Regiane Souza Camargo Moreira apresentou o trabalho “¿ Cómo evaluar el impacto de un rompe-hielo en proyecto COIL?” (Como avaliar o impacto de um quebra-gelo no projeto COIL?), elaborado em coautoria com Javier Guerrero e Sandra Cruz (Uniminuto / Colômbia).
Destacou que o quebra-gelo promove conexão e diálogo intercultural entre os participantes, daí a importância de seu planejamento.
É essencial se preparar para poder se expor, especialmente nos aspectos pessoais e culturais. “O quebra-gelo não garante o \textbf{sucesso} do projeto, mas ajuda a reforçar o trabalho em equipe e a evolução do COIL”, afirmou Regiane, que compartilhou pesquisa realizada com 30 professores e 4 coordenadores COIL.
Ana Carolina Freschi fez duas comunicações relacionadas à \textbf{institucionalização} dos projetos COIL.
A primeira foi “GSL Classroom and the Chamber of Secrets: The Journey throughout Institutionalizations of COIL Initiatives” (Intercâmbios Virtuais e a Câmara Secreta: a \textbf{jornada} pelas \textbf{institucionalizações} das iniciativas COIL) com Ricardo Lyle Bañuelos, Gisselle Morales e Monserrat Balandrano (Tecnológico de Monterrey / México), Dan Nolan (University of Minnesota System / EUA) e Natalia Dyba (University of Washington Bothell / EUA).
Os apresentadores destacaram a importância da melhoria contínua dos processos de \textbf{institucionalização}, com diretrizes e requisitos detalhados.
O Tecnológico de Monterrey, por exemplo, trabalha com chamada de propostas para conhecer os planos de colaboração internacional dos professores interessados em iniciativas GSL (Global Shared Learning, que é como a instituição chama os Intercâmbios Virtuais).

% matched lemmas: institucionalização, jornada, sucesso
\end{textsample}
